\section{Personal Revelations}

This post will address some of the more common questions that readers have asked me.

\paragraph{Do you Blog full time?}


With the massive amount of all original content, new translations, not to mention the distasteful duties of a webmaster, maintaining a full-featured web site may look like a full time job. But, no, it is a part time affair on top of the normal duties of holding a job, running a household, and so on. However, my energy level is decreasing and it needs to be a full time job, or else it will terminate. There are associated expenses to pay for web hosting and research materials from Europe. Please consider responding to the poll questions at the end for suggestions about defraying expenses.

There is a convention among Hermetists to not earn a living through the Work. Back in the day, the Hermetist would have held a job like a horse trader, or a street artist that would give him cover as he traveled from town to town on market day or other festivals. Since books were expensive and bulky, he would carry around the teachings encoded in diagrams or mandalas. These were portable and could be copied by the students. Eventually, some of them were incorporated into the Tarot deck. Along with the oral teaching, they served as a condensed digest of the work until the teacher's next visit.

\paragraph{What does Gornahoor mean?}


Gornahoor is simply a title of respect like sir, mister, or sri, although it is reserved for a three-brained being.

\paragraph{Are you Eastern Orthodox?}


Several people have assumed this. Obviously, for more than 1000 years Rome was Orthodox and Orthodoxy was Roman. So I don't recognize the schism, as it is the period prior to the Reformation and the Renaissance in the West that is of the most interest. Furthermore, I would (figuratively) agree with this quote, at a time when Constantinople, rather than Athens, was the center of the East:

\begin{quotex}
Constantinople remains to this day a monument of ancient wisdom, the home of letters, and the citadel of philosophy. No Italian can pretend to be educated unless he has studied for a time at Constantinople.
\flright{\textsc{Aeneas Sylvius (1551)}}
\end{quotex}


\paragraph{What religion do you follow?}


With few exceptions, a man ought to follow the exoteric religion of his nation, as it is the creation of its ethnic soul. People, even those who claim to be close to ``tradition'', today treat their religion as a commodity, all sitting on a shelf, with the consumer choosing one based on his personal peccadilloes. They may complain because they don't like what the produces has provided them. That is the passive attitude. Once it is understood that it is a reflection of the spirit of a people, the choices looks much different. I don't have time to go into this now, so I am reserving it to some future discussions on the idealistic philosophy of Giovanni Gentile. It may make things comprehensible.

What is called ``Traditional Catholicism'', or Rad Trad, is not the same thing. Rad Trad really goes back to the Council of Trent, which is hampered by its need to address the false claims of the Reformation. Thus, too, its spirituality is primarily influenced by the post-Tridentine saints from Spain and France. Not that this is bad in itself, but it is necessary to remain further back in tradition.

Despite the prevalence of Guenon's works in our time, at least in certain circles, it is astounding how poorly understood he still is. Exoterism and esoterism are of different orders and, as such, are incomparable. This is lost on a lot of readers who attempt to find a connection or correspondence between exoteric dogmas and esoteric teachings. This shows up in two ways. The first are those orthodox believers who are concerned that some dogma may be transgressed. The worst ones are those who try to ``prove'' that Christianity, which for us can only mean the Roman/Orthodox synthesis, is not ``traditional'' or ``esoteric''. Of course, the exoteric is not esoteric, and the latter has never been fully revealed to the public at large, who would either misunderstand it or ignore it altogether.

As Guenon reassures us:

\begin{quotex}
A truly initiatic tradition cannot be heterodox; to quality it as such is to reverse the normal and hierarchical relationship between the interior and the exterior. Esoterism is not contrary to orthodoxy, even orthodoxy construed simply in the religious sense; it is above or beyond the religious point of view, which is obviously not at all the same thing.
\end{quotex}


That there have been conflicts between the two perspectives in the past has been due to completely other factors. More recently, Valentin Tomberg is firm in his insistence that he has no intention to challenge any established dogmas. I concur with that.

\paragraph{How can a university graduate believe such things?}


There is more than a millennium of very sophisticated thought behind the teachings, reaching up to the stunning synthesis of Thomas Aquinas. As Thomas Aquinas himself admitted, the way he understands church teaching would be quite different from that of the typical believer.

As he points out, dogma demands the ``free consent of the will''. Obviously, only a free man can grant such consent. A dogma is an internal constraint, not external, and it is a point of view. The bulk of humanity chooses their opinions and points of view with very little self-reflection; dogmas, on the other hand, required much reflection, detachment, and loyalty to one's tradition.

Most men would object: can you prove it is true (like the Greeks), or can you give me a sign (like the Jews). But then the choice would not be free, it would be compelled. Think about this a little bit and you will see it is true. Furthermore, few things can really be proven true and signs can be interpreted in opposite ways.

Since the free man transcends all points of view, he can choose the one most appropriate for him.

\paragraph{Why the Fedeli D'Amore?}


Obviously because it provides an open example of an esoteric current. Unlike a Reghini or a Rossetti, Guenon points out that the Fedeli d'Amore were not heterodox for the reasons given. An intriguing point is the role that Medusa played in their symbolism. Now Dante was initiated by the Sicilian troubadours and Medusa to this day is on the flag of Sicily. A couple of years ago, a group tried to change it to something less pagan, but they were rebuffed.

Without further explanation, catholic means universal, the one and only tradition. It is the new Israel, i.e., the totality of initiates, or the chosen people. In Guenon's words,

Has not the name Israel been employed to designate the totality of initiates and who possess the universal language that enables them all to understand each other, that is, the knowledge of the one tradition that is concealed beneath all its particular forms?

That is exactly what Tomberg says about the method of depth. Hence, Dante really means the Church is teaching the universal tradition and he is not opposing it as such. Those who see things differently, are most likely restricted to the horizontal, profane plane.

\paragraph{What's love got to do with it?}


Amor is Roma spelled backwards and may be related to the occult name of Rome. Some, unfortunately, see A-mor as anti-Rome rather than as its mirror reflection.

When an esoteric writing does not make much sense, that is probably because you are missing the deeper meaning. Clearly, Dante's crush on Beatrice exceeds all normal bounds. Yet it is that awakening of eros that provides the energy and impetus for the spiritual climb. That might not help you, if you have never experienced a real love for a woman, or are not even capable of imagining such a thing.

The novice to the Fedeli d'Amore had to write a love sonnet. That would be a good homework assignment, especially if you could sublimate that feeling to the Divine Sophia. Remember this is quite a different feeling from that associated with watching Internet porn. The least serious sexual sin is intercourse between an unmarried man and woman; solitary pornography does not satisfy that and can only result in unnatural acts. There are sound reasons for this.

However, don't think that accumulating sexual experiences is an acceptable substitute. When you get old, you remember none of them. Instead, you become obsessed with the ones you let get away, or the acts you never got to try. It is far better to be obsessed with the rebis to achieve the spiritual marriage which is the primordial state.

The echoes of that tradition of Amore can be found in the tradition of Neapolitan love songs, although much secularized. However, is you need help capturing that feeling, try listening to
\emph{Core N'grato}\footnote{\url{https://www.youtube.com/watch?v=NOVBrZzjOkQ} by \textbf{Giuseppe di Stefano}}. Unfortunately, such emotions are often wasted on a woman, who only gets so passionate when shopping.

\textbf{How can Gornahoor generate income}


\begin{itemize}


\item Find a wealthy benefactor {(25\%, 51 Votes)}

\item Beg for donations {(20\%, 42 Votes)}

\item Sell monogrammed t-shirts, coffee mugs, lapel pins, busts of famous traditionalists, etc. {(16\%, 33 Votes)}

\item Charge for weekend seminars to determine your spiritual race and purpose in life. {(16\%, 32 Votes)}

\item Produce a reality TV show called the "Boys of Tradition" (15\%, 30 Votes)

\item Make Gornahoor a subscription only web site {(5\%, 11 Votes)}

\item Other (specify in comments) {(3\%, 7 Votes)}
\end{itemize}


Total Voters: \textbf{118}

\url{https://www.youtube.com/watch?v=NOVBrZzjOkQ}

\flrightit{Posted on 2013-04-08 by Cologero}

\begin{center}* * *\end{center}

\begin{footnotesize}
\begin{sffamily}
\texttt{Nick on 2013-04-09 at 03:15 said:}

What would the Templars do?

\hfill

\texttt{Mihai on 2013-04-09 at 03:39 said:}

Your polls always make me burst into laughter, Cologero.

\hfill

\texttt{Ea Lord of the Deeps on 2013-04-09 at 07:27 said:}

I would buy a Evola related tshirt,

\hfill

\texttt{Cavalcare on 2013-04-09 at 08:33 said:}

Laptop stickers. An example would be a certain picture circulating around the tumblr-sphere: ``Uber'', picture of Nietzsche above, ``Decline'', picture of Spengler, ``Tradition'', picture of Evola. Can't seem to track it down.

\hfill

\texttt{Avery Morrow on 2013-04-09 at 09:39 said:}

I would attend a Gornahoor weekend retreat, just to see if any of us could transcend our cynical detachment.

\hfill

\texttt{J. Alexander Maximilian on 2013-04-09 at 09:58 said:}

I think donations would be good. Id be happy to give what little I could. Or subscribe at \$10 a month. Or ``things'' like t-shirts, etc.

\hfill

\texttt{J. Alexander Maximilian on 2013-04-09 at 10:08 said:}

And I have loved a woman with much Romance, but the energies invested therein are undoubtedly ``wasted'' as you said.

The misery of Amoré at times. The heart JT leaps, and is quickly shot down eventually by the shopper.

\hfill

\texttt{Logres on 2013-04-09 at 11:06 said:}

I would be happy to contribute what money I can send, although I am not sure everyone that needs to read it is able to give anything. Or, you could take donations for certain ends, like meeting website cost or getting a new manuscript from Europe.

\hfill

\texttt{Jason-Adam on 2013-04-09 at 12:43 said:}

I personally would be like Evola or Yockey and use my noble title to seduce a rich older lady with my virile masculinity who would then pay my bills to support my Traditionalist work which she could never understand as she is an inferior woman but would remain unable to resist the passion it gives her in my lovemaking... ... ..this comment is a joke. Or is it ?

\hfill

\texttt{Anna Kaiser-Swadling on 2013-04-09 at 14:18 said:}

You could replace the Tao of shopping.

\hfill

\texttt{Matt on 2013-04-09 at 20:47 said:}

Honestly, I wouldn't be opposed to paying a reasonable subscription price to Gornahoor. I don't think many people would be all that interested in seeing a show with people who take up camera time searching for truth and cosmic redemption, and who also don't look like carrots after days of over-tanning.

To echo Avery's comment, I also think seminar retreats is a good idea. Gives us a chance to meet each other and Cologero in person, and a chance for mutual support and guidance.

\hfill

\texttt{David on 2013-04-09 at 22:05 said:}

If I can get the whole bunch of all articles in a .rar or .zip files, I would be glad to pay for it. That would be most useful.

\hfill

\texttt{George Goerlich on 2013-04-10 at 00:16 said:}

Can you sell your translations?

\hfill

\texttt{Mihai on 2013-04-10 at 03:33 said:}

One problem: many of us are literally continents apart. A plane ticket from one continent to another isn't all that accessible nowadays.

That is, unless, we arrange a non-corporeal meeting at the primordial center of the present aeon, preferably by the shade provided by the Tree in the middle. But they told me they put a Cherub at the entrance and it may take even harder work to get passed his flaming sword than the option with the plane ticket.

\hfill

\texttt{VerbaVolant on 2013-04-10 at 13:23 said:}

thanks for quoting the great René Guénon !!

\hfill

\texttt{VerbaVolant on 2013-04-10 at 13:25 said:}

René Guénon is surely \& definetely a great reference !!

\hfill

\texttt{Ammonius420 on 2013-04-10 at 20:15 said:}

Apparently not:

\url{http://www.youtube.com/watch?v=MOdb1TAkh3k}


``He also dated a woman around 20-30 years his senior too, I was friends with him on Facebook before.? Quite an unconventional character to mildly say :o''

P.S. I would pay for this site.

\hfill

\texttt{Ammonius420 on 2013-04-11 at 03:32 said:}

How about finding ways to increase Cologero's energy level?

\hfill

\texttt{Cologero on 2013-04-11 at 08:45 said:}

From Wikepedia article on the Hermetist, Giacomo Casanova:

\begin{quotationx}

Escaping to Parma, Casanova entered into a three-month affair with a Frenchwoman he named ``Henriette'', perhaps the deepest love he ever experienced---a woman who combined beauty, intelligence, and culture. In his words, ``They who believe that a woman is incapable of making a man equally happy all the twenty-four hours of the day have never known an Henriette. The joy which flooded my soul was far greater when I conversed with her during the day than when I held her in my arms at night. Having read a great deal and having natural taste, Henriette judged rightly of everything.'' She also judged Casanova astutely. As noted Casanovist J. Rives Childs wrote:

Perhaps no woman so captivated Casanova as Henriette; few women obtained so deep an understanding of him. She penetrated his outward shell early in their relationship, resisting the temptation to unite her destiny with his. She came to discern his volatile nature, his lack of social background, and the precariousness of his finances. Before leaving, she slipped into his pocket five hundred louis, mark of her evaluation of him.
\end{quotationx}


I don't know that 500 louis are all that great for three months work, not to mention he also gave up his heart to her. I'm afraid my peasant background will forever bar me from such circles, despite my not unsubstantial charms.

\hfill

\texttt{Cologero on 2013-04-11 at 09:05 said:}

After the alchemist Fulcanelli wrote \textit{The Mystery of the Cathedrals} --- which did for the understanding of the symbolism of the Medieval cathedrals what Luigi Valli did for understanding the symbolism of the Fedeli d'Amore --- he disappeared without a trace. When he did arrange a clandestine meeting with a pupil several years later, the pupil reported that Fulcanelli has actually become younger. I know what I need to do.

\hfill

\texttt{Jason-Adam on 2013-04-11 at 13:12 said:}

Ammonius, who is that weirdo communist in that video you posted ? I am not sure what you intended by choosing someone with the same first name as I but he is funny though he looks a bit dazed as if he were using some illegal substance... ..

and as for Cologero, you are a peasant ? You can not be, your spirit as I read it in your writings here proves you are an aristocrat... ... ... 

\hfill

\texttt{Ammonius420 on 2013-04-11 at 19:39 said:}

I wish you luck.

\hfill

\texttt{CaseyAnn on 2013-04-21 at 18:56 said:}

But aren't t-shirts a hideous product of the modern age . . . ?

:-p

\hfill

\texttt{Cologero on 2013-04-21 at 21:05 said:}

I'll take that as a ``no'' vote. Perhaps you can choose something from \textit{this list}\footnote{\url{https://www.google.com/search?q=medieval+clothes\&hl=en\&tbm=isch\&tbo=u\&source=univ\&sa=X\&ei=q4x0Ubr9PIq09QSe5oG4Dg\&sqi=2\&ved=0CD8QsAQ\&biw=1009\&bih=568}}?

\hfill

\texttt{CaseyAnn on 2013-04-21 at 22:25 said:}

Ha! But that's not a cynical dichotomy, is it?

I voted awhile ago for the now bolded, italicized text.

Question about the Catholic/Orthodox synthesis: How would that look in practice? Exoterically, I mean. I find myself drawn to this but don't understand.

I really enjoy this site, Cologero. Thank you.

\hfill

\texttt{Cologero on 2013-04-21 at 23:58 said:}

I don't recall having mentioned any such ``Catholic/Orthodox synthesis''. What is there to synthesize? The West regards the separation as a ``schism'', i.e., a dispute over jurisdiction, not theology. I know the Easterners are always quick to identify a heresy where there really isn't one. The break up almost turned into a makeup at the Council of Florence a few hundred years ago, but extraneous issues intervened. In any event, I believe the mutual ex-communications were lifted in the 1960s, although, for some reason that makes no sense, there is still no full communion.

In practice, I suggest studying the Philokalia. You may find it simpatico. Or else go to a Greek festival where you may find an icon of St. Augustine, as I did.

\hfill

\texttt{Jason-Adam on 2013-04-22 at 12:13 said:}

The problem is geopolitical -- the East believes the West should be subjected to Constantinople (in the middle ages) and Russia (today -- read the works of Alexander Dugin) which causes the eastern church to ally with neopagans against the Western church.

We are westerners who want the west to return to tradition of our ancestors, the east does not care about our tradition, they want to destroy us -- read the debate between Olavo and Dugin and you will see that Dugin does not hate the modern world -- he hates the West in toto.

The Photian ``orthodox'' heretics believe that the Holy Spirit does not proceed from the Son but He comes independently from the Father alone -- this leads to the idea of multiple voices of God and hence the principle of chaos being the ruling Idea of the world.

\hfill

\texttt{Boreas on 2020-04-08 at 10:41 said:}

What happened to the donate button?! I'd very much like to support your work with the little I have.

\hfill

\texttt{DA on 2022-02-26 at 18:08 said:}

How about self-publishing the collections of articles available here, organized in certain topics? I'd definitely buy that, paying a higher margin to cover the expenses that'll keep Gornahoor running.

Also a Patreon-like monthly donations are definitely part of the equation for me. You could attach perks to certain levels of financial support. Eg. 5 USD/month -- your name mentioned in the list of benefactors. 10 USD/month -- previous benefit + one extra access article. 50 USD/month -- all the former and a Zoom coffee meeting with the author? Treat numbers as arbitrary, I'm your reader from Poland and I don't know what's considered appropriate for such initiatives in the US.

All the best,

Heiko

\hfill

\texttt{Cologero on 2022-03-10 at 22:16 said:}

Actually, DA, I am looking for some software that will take a wordpress database and convert it to something like a Word doc or PDF file. Then I could re-organize the articles by topic. If anyone knows of such software, please let me know.

In the USA, anything that ends with USD is appropriate.

\end{sffamily}
\end{footnotesize}

